\documentclass[10pt,a4paper]{amsart}
\usepackage{amsmath,amssymb,amsthm,graphicx}
% Retained preamble for the Stacks Project chapter morphisms.tex (Morphisms of Schemes).
% The chapter file pinned-original-source.tex carries only \input{preamble} on its first line, so no
% preamble text lives inside the chapter itself.  The Stacks Project uses one single shared preamble for
% every chapter; the snapshot stages it verbatim in sources/stacks-chow/preamble.tex.  The block below is
% copied mechanically, line for line, from that file's lines 14-18 (the Xy-pic setup) and lines 154-217
% (the native theorem declarations together with the mathematical macros).  Nothing was added, removed or
% redefined.  The wrapper class and the amsmath/amssymb/amsthm/graphicx base packages are supplied by the
% assembler, so the class/\externaldocument part of the shared preamble is not repeated here.
% For commutative diagrams we use Xy-pic
\usepackage[all]{xy}

% We use 2cell for 2-commutative diagrams.
\xyoption{2cell}
% Theorem environments.
%
\theoremstyle{plain}
\newtheorem{theorem}[subsection]{Theorem}
\newtheorem{proposition}[subsection]{Proposition}
\newtheorem{lemma}[subsection]{Lemma}

\theoremstyle{definition}
\newtheorem{definition}[subsection]{Definition}
\newtheorem{example}[subsection]{Example}
\newtheorem{exercise}[subsection]{Exercise}
\newtheorem{situation}[subsection]{Situation}

\theoremstyle{remark}
\newtheorem{remark}[subsection]{Remark}
\newtheorem{remarks}[subsection]{Remarks}

\numberwithin{equation}{subsection}

% Macros
%
\def\lim{\mathop{\mathrm{lim}}\nolimits}
\def\colim{\mathop{\mathrm{colim}}\nolimits}
\def\Spec{\mathop{\mathrm{Spec}}}
\def\Hom{\mathop{\mathrm{Hom}}\nolimits}
\def\Ext{\mathop{\mathrm{Ext}}\nolimits}
\def\SheafHom{\mathop{\mathcal{H}\!\mathit{om}}\nolimits}
\def\SheafExt{\mathop{\mathcal{E}\!\mathit{xt}}\nolimits}
\def\Sch{\mathit{Sch}}
\def\Mor{\mathop{\mathrm{Mor}}\nolimits}
\def\Ob{\mathop{\mathrm{Ob}}\nolimits}
\def\Sh{\mathop{\mathit{Sh}}\nolimits}
\def\NL{\mathop{N\!L}\nolimits}
\def\CH{\mathop{\mathrm{CH}}\nolimits}
\def\proetale{{pro\text{-}\acute{e}tale}}
\def\etale{{\acute{e}tale}}
\def\QCoh{\mathit{QCoh}}
\def\Ker{\mathop{\mathrm{Ker}}}
\def\Im{\mathop{\mathrm{Im}}}
\def\Coker{\mathop{\mathrm{Coker}}}
\def\Coim{\mathop{\mathrm{Coim}}}

% Boxtimes
%
\DeclareMathSymbol{\boxtimes}{\mathbin}{AMSa}{"02}

%
% Macros for moduli stacks/spaces
%
\def\QCohstack{\mathcal{QC}\!\mathit{oh}}
\def\Cohstack{\mathcal{C}\!\mathit{oh}}
\def\Spacesstack{\mathcal{S}\!\mathit{paces}}
\def\Quotfunctor{\mathrm{Quot}}
\def\Hilbfunctor{\mathrm{Hilb}}
\def\Curvesstack{\mathcal{C}\!\mathit{urves}}
\def\Polarizedstack{\mathcal{P}\!\mathit{olarized}}
\def\Complexesstack{\mathcal{C}\!\mathit{omplexes}}
% \Pic is the operator that assigns to X its picard group, usage \Pic(X)
% \Picardstack_{X/B} denotes the Picard stack of X over B
% \Picardfunctor_{X/B} denotes the Picard functor of X over B
\def\Pic{\mathop{\mathrm{Pic}}\nolimits}
\def\Picardstack{\mathcal{P}\!\mathit{ic}}
\def\Picardfunctor{\mathrm{Pic}}
\def\Deformationcategory{\mathcal{D}\!\mathit{ef}}
\usepackage{cleveref}
\crefname{originalreference}{original reference}{original references}
\Crefname{originalreference}{Original reference}{Original references}
\makeatletter
\newlabel{algebra-lemma-differentials-diagonal}{{algebra-lemma-differentials-diagonal}{source}{}{sourcepointer-1}{}}
\newlabel{algebra-lemma-differentials-diagonal@cref}{{[originalreference][0][]algebra-lemma-differentials-diagonal}{[1][0][]source}{}{}{}}
\makeatother
\crefname{originaldefinition}{original definition}{original definitions}
\Crefname{originaldefinition}{Original definition}{Original definitions}
\makeatletter
\newlabel{definition-conormal-sheaf}{{definition-conormal-sheaf}{source}{}{sourcepointer-2}{}}
\newlabel{definition-conormal-sheaf@cref}{{[originaldefinition][0][]definition-conormal-sheaf}{[1][0][]source}{}{}{}}
\makeatother
\crefname{originallemma}{original lemma}{original lemmas}
\Crefname{originallemma}{Original lemma}{Original lemmas}
\makeatletter
\newlabel{lemma-affine-conormal}{{lemma-affine-conormal}{source}{}{sourcepointer-3}{}}
\newlabel{lemma-affine-conormal@cref}{{[originallemma][0][]lemma-affine-conormal}{[1][0][]source}{}{}{}}
\makeatother
\crefname{originallemma}{original lemma}{original lemmas}
\Crefname{originallemma}{Original lemma}{Original lemmas}
\makeatletter
\newlabel{lemma-differentials-affine}{{lemma-differentials-affine}{source}{}{sourcepointer-4}{}}
\newlabel{lemma-differentials-affine@cref}{{[originallemma][0][]lemma-differentials-affine}{[1][0][]source}{}{}{}}
\makeatother
\crefname{originalreference}{original reference}{original references}
\Crefname{originalreference}{Original reference}{Original references}
\makeatletter
\newlabel{modules-lemma-double-structure-gives-derivation}{{modules-lemma-double-structure-gives-derivation}{source}{}{sourcepointer-5}{}}
\newlabel{modules-lemma-double-structure-gives-derivation@cref}{{[originalreference][0][]modules-lemma-double-structure-gives-derivation}{[1][0][]source}{}{}{}}
\makeatother
\crefname{originalreference}{original reference}{original references}
\Crefname{originalreference}{Original reference}{Original references}
\makeatletter
\newlabel{schemes-lemma-diagonal-immersion}{{schemes-lemma-diagonal-immersion}{source}{}{sourcepointer-6}{}}
\newlabel{schemes-lemma-diagonal-immersion@cref}{{[originalreference][0][]schemes-lemma-diagonal-immersion}{[1][0][]source}{}{}{}}
\makeatother
\title{The sheaf of relative differentials is canonically isomorphic to the conormal sheaf of the diagonal}
\author{The Stacks Project Authors}
\date{stacks-project:morphisms.tex:local-snapshot-20261012; GNU FDL 1.2 or later}
\begin{document}
\maketitle
\noindent\textit{Editorial note.} All delivered mathematical text is the author's own, copied from the pinned Stacks Project chapter file pinned-original-source.tex; no statement or proof step was rewritten, reordered or removed.\par\medskip
\noindent\textit{Editorial note.} Following the multi-claim convention for this source, source\_url is the source-level URL https://stacks.math.columbia.edu/ and the source label lemma-differentials-diagonal is carried in claim\_key and here rather than in a per-claim tag URL.\par\medskip
\noindent\textit{Editorial note.} The preamble is the chapter's shared Stacks Project preamble, retained as sources/stacks-morphisms/retained-preamble.tex: a line-for-line excerpt of the shared Stacks preamble staged with this snapshot. No macro was invented or redefined, and this card adds no preamble text of its own.\par\medskip
\noindent\textit{Editorial note.} Statement span 6506-6511 and proof span 6513-6595 come from the halfA pool card operations/pool-cards/stacks-morphisms-halfA.json and were re-read in the pinned chapter file before this card was written.\par\medskip
\noindent\textit{Editorial note.} The proof cites two labels of the same chapter (the conormal sheaf definition and the affine conormal lemma) and several labels of other chapters (Schemes, Modules, Algebra); all of them are delivered as external navigation pointers only.\par\medskip
\noindent\textit{Editorial note.} 6 cross-references occur inside the counted spans; they are listed in external\_references as navigation pointers only.\par\medskip
\noindent\textit{Allowed standard prerequisites (editorial).}\begin{itemize}
\item The diagonal morphism is an immersion, so that it has an ideal sheaf on an open subscheme of the fibre product containing its image; the author invokes the same-chapter statement for this and for the shape of the chosen open. Invoked at proof lines 6515-6519 and again at 6551-6554. Supplies the open W and the ideal sheaf whose square-zero thickening carries the two sections.
\item The conormal sheaf of an immersion, its definition, and its description over an affine chart as I/I\textasciicircum{}2, together with the affine description of the module of differentials. Definition and lemmas of the same chapter cited at proof lines 6526-6528 and 6572-6576. Names the target of the canonical map and lets the affine computation identify both sides with J/J\textasciicircum{}2 and with the module of differentials of the coordinate ring.
\item The double-structure mechanism: a ring surjection A to B carrying two sections gives a derivation B to the kernel of the surjection, and the universal property of the module of differentials turns it into a linear map. Cited as an external chapter lemma at proof lines 6530 and 6540, and applied to the square-zero thickening O\_W/J\textasciicircum{}2 at proof lines 6540-6548. Produces the canonical map from the module of differentials to the conormal sheaf.
\end{itemize}\medskip
\ifdefined\hypertarget\hypertarget{sourcepointer-1}{}\fi
\noindent\textit{Source reference.} Original reference algebra-lemma-differentials-diagonal: sources/stacks-morphisms/pinned-original-source.tex, source lines 6594--6594. Same-chapter reference used in the retained author text at source line(s) 6594. It is supplied only as a navigation pointer so that the retained text compiles; the counted item delivers the author's own statement and the author's complete proof, and no part of the referenced item is copied in or substituted for the counted argument.\par\medskip
\ifdefined\hypertarget\hypertarget{sourcepointer-2}{}\fi
\noindent\textit{Source reference.} Original definition definition-conormal-sheaf: sources/stacks-morphisms/pinned-original-source.tex, source lines 6164--6164. Same-chapter reference used in the retained author text at source line(s) 6537. It is supplied only as a navigation pointer so that the retained text compiles; the counted item delivers the author's own statement and the author's complete proof, and no part of the referenced item is copied in or substituted for the counted argument.\par\medskip
\ifdefined\hypertarget\hypertarget{sourcepointer-3}{}\fi
\noindent\textit{Source reference.} Original lemma lemma-affine-conormal: sources/stacks-morphisms/pinned-original-source.tex, source lines 6186--6186. Same-chapter reference used in the retained author text at source line(s) 6538, 6588. It is supplied only as a navigation pointer so that the retained text compiles; the counted item delivers the author's own statement and the author's complete proof, and no part of the referenced item is copied in or substituted for the counted argument.\par\medskip
\ifdefined\hypertarget\hypertarget{sourcepointer-4}{}\fi
\noindent\textit{Source reference.} Original lemma lemma-differentials-affine: sources/stacks-morphisms/pinned-original-source.tex, source lines 6441--6441. Same-chapter reference used in the retained author text at source line(s) 6495, 6590, 6701, 6716, 7660. It is supplied only as a navigation pointer so that the retained text compiles; the counted item delivers the author's own statement and the author's complete proof, and no part of the referenced item is copied in or substituted for the counted argument.\par\medskip
\ifdefined\hypertarget\hypertarget{sourcepointer-5}{}\fi
\noindent\textit{Source reference.} Original reference modules-lemma-double-structure-gives-derivation: sources/stacks-morphisms/pinned-original-source.tex, source lines 6546--6546. Same-chapter reference used in the retained author text at source line(s) 6546, 6559. It is supplied only as a navigation pointer so that the retained text compiles; the counted item delivers the author's own statement and the author's complete proof, and no part of the referenced item is copied in or substituted for the counted argument.\par\medskip
\ifdefined\hypertarget\hypertarget{sourcepointer-6}{}\fi
\noindent\textit{Source reference.} Original reference schemes-lemma-diagonal-immersion: sources/stacks-morphisms/pinned-original-source.tex, source lines 6520--6520. Same-chapter reference used in the retained author text at source line(s) 6520, 6525, 6571. It is supplied only as a navigation pointer so that the retained text compiles; the counted item delivers the author's own statement and the author's complete proof, and no part of the referenced item is copied in or substituted for the counted argument.\par\medskip
\begin{lemma}
\label{lemma-differentials-diagonal}
Let $f : X \to S$ be a morphism of schemes. There is a canonical
isomorphism between $\Omega_{X/S}$ and the conormal sheaf of
the diagonal morphism $\Delta_{X/S} : X \longrightarrow X \times_S X$.
\end{lemma}

\begin{proof}
We first establish the existence of a couple of ``global'' sheaves
and global maps of sheaves, and further down we describe
the constructions over some affine opens.

\medskip\noindent
Recall that $\Delta = \Delta_{X/S} : X \to X \times_S X$
is an immersion, see Schemes, Lemma \ref{schemes-lemma-diagonal-immersion}.
Let $\mathcal{J}$ be the ideal sheaf of the immersion
which lives over some open subscheme $W$ of $X \times_S X$
such that $\Delta(X) \subset W$ is closed. Let us take the one that
was found in the proof of
Schemes, Lemma \ref{schemes-lemma-diagonal-immersion}.
Note that the sheaf of rings $\mathcal{O}_W/\mathcal{J}^2$
is supported on $\Delta(X)$. Moreover it sits in a
short exact sequence of sheaves
$$
0 \to \mathcal{J}/\mathcal{J}^2
\to \mathcal{O}_W/\mathcal{J}^2
\to \Delta_*\mathcal{O}_X
\to 0.
$$
Using $\Delta^{-1}$ we can think of this as a surjection of
sheaves of $f^{-1}\mathcal{O}_S$-algebras with kernel the
conormal sheaf of $\Delta$ (see Definition \ref{definition-conormal-sheaf}
and Lemma \ref{lemma-affine-conormal}).
$$
0 \to \mathcal{C}_{X/X \times_S X}
\to \Delta^{-1}(\mathcal{O}_W/\mathcal{J}^2)
\to \mathcal{O}_X
\to 0
$$
This places us in the situation of
Modules, Lemma \ref{modules-lemma-double-structure-gives-derivation}.
The projection morphisms $p_i : X \times_S X \to X$, $i = 1, 2$ induce
maps of sheaves of rings
$(p_i)^\sharp : (p_i)^{-1}\mathcal{O}_X \to \mathcal{O}_{X \times_S X}$.
We may restrict to $W$ and quotient by $\mathcal{J}^2$ to get
$(p_i)^{-1}\mathcal{O}_X \to \mathcal{O}_W/\mathcal{J}^2$.
Since $\Delta^{-1}p_i^{-1}\mathcal{O}_X = \mathcal{O}_X$
we get maps
$$
s_i : \mathcal{O}_X \to \Delta^{-1}(\mathcal{O}_W/\mathcal{J}^2).
$$
Both $s_1$ and $s_2$ are sections to the map
$\Delta^{-1}(\mathcal{O}_W/\mathcal{J}^2) \to \mathcal{O}_X$,
as in Modules, Lemma \ref{modules-lemma-double-structure-gives-derivation}.
Thus we get an $S$-derivation
$\text{d} = s_2 - s_1 : \mathcal{O}_X \to \mathcal{C}_{X/X \times_S X}$.
By the universal property of the module of differentials we find a
unique $\mathcal{O}_X$-linear map
$$
\Omega_{X/S} \longrightarrow \mathcal{C}_{X/X \times_S X},\quad
f\text{d}g \longmapsto fs_2(g) - fs_1(g)
$$
To see the map is an isomorphism, let us work this out over suitable affine
opens. We can cover $X$ by affine opens $\Spec(A) = U \subset X$
whose image is contained in an affine open $\Spec(R) = V \subset S$.
According to the proof of Schemes, Lemma \ref{schemes-lemma-diagonal-immersion}
$U \times_V U \subset X \times_S X$ is an affine open
contained in the open $W$ mentioned above. Also
$U \times_V U = \Spec(A \otimes_R A)$.
The sheaf $\mathcal{J}$ corresponds to the ideal
$J = \Ker(A \otimes_R A \to A)$.
The short exact sequence to the short exact sequence
of $A \otimes_R A$-modules
$$
0 \to J/J^2 \to (A \otimes_R A)/J^2 \to A \to 0
$$
The sections $s_i$ correspond to the ring maps
$$
A \longrightarrow (A \otimes_R A)/J^2,\quad
s_1 : a \mapsto a \otimes 1,\quad
s_2 : a \mapsto 1 \otimes a.
$$
By Lemma \ref{lemma-affine-conormal} we have
$\Gamma(U, \mathcal{C}_{X/X \times_S X}) = J/J^2$ and
by Lemma \ref{lemma-differentials-affine}
we have $\Gamma(U, \Omega_{X/S}) = \Omega_{A/R}$.
The map above is the map $a \text{d}b \mapsto a \otimes b - ab \otimes 1$
which is shown to be an isomorphism in
Algebra, Lemma \ref{algebra-lemma-differentials-diagonal}.
\end{proof}

\end{document}
